\documentclass[11pt]{article}
\usepackage[margin=1in]{geometry}
\usepackage{amsmath,amssymb,mathtools,bm}

\title{Janiform Compact Model Spec}
\author{}
\date{}

\begin{document}
    \maketitle

    \section*{1. Return Process}
    \[
        X_t^{\mathrm{JF}} = L^+_{A_t^+} - L^-_{A_t^-} + M_t + G_t,
        \qquad
        A_t^\pm = \int_0^t \lambda_u^\pm\,du.
    \]

    \section*{2. Observed-State Mean Reversion Layer}
    Let \(\Sigma_h\) denote the structural Janiform boundary state on interval \(h\), and let
    \[
        K_h^P(u \mid \Sigma_h)
        =
        \log \mathbb E^P[e^{u R_h}\mid \Sigma_h]
    \]
    be the one-interval physical log-moment-generating function of the Janiform core contribution \(R_h\).

    Define the physical expected core increment
    \[
        \mu_h^P(\Sigma_h)
        =
        \partial_u K_h^P(u \mid \Sigma_h)\big|_{u=0}
        =
        \mathbb E^P[R_h \mid \Sigma_h].
    \]

    Let \(\Delta_{\mathrm{ref}} > 0\) be the persisted reference interval, let
    \(\psi_{\mathrm{ref}} \in [0,1]\) be the persisted pressure persistence per
    reference interval, let \(c \in [0,1]\) be the persisted pressure-to-level
    gain, and define
    \[
        \psi_h = \psi_{\mathrm{ref}}^{\Delta_h / \Delta_{\mathrm{ref}}}.
    \]

    Let the observed state satisfy
    \[
        X_h = A_h^P + D_h^P,
        \qquad
        A_h^P = a_0^P + c\,p_h^P,
        \qquad
        p_0^P = 0,
        \qquad
        \phi_h=e^{-\kappa \Delta_h},
        \qquad
        \kappa \ge 0.
    \]

    Then the endogenous pressure-to-level recursions are
    \begin{gather*}
        p_{h+1}^P = \psi_h p_h^P + (1-\psi_h)\mu_h^P(\Sigma_h),\\
        A_{h+1}^P = A_h^P + c\big(p_{h+1}^P - p_h^P\big),\\
        D_{h+1}^P
        =
        \phi_h D_h^P + R_h + A_h^P - A_{h+1}^P
        =
        \phi_h D_h^P + R_h - c\big(p_{h+1}^P - p_h^P\big),\\
    \end{gather*}
    and therefore
    \[
        X_{h+1}
        =
        A_{h+1}^P + D_{h+1}^P
        =
        X_h + (\phi_h - 1)D_h^P + R_h.
    \]

    \section*{3. Trade Intensities and Hawkes State}
    \[
        \lambda_t^+ = s(t)\bar{\lambda} + Y_t^{++} + Y_t^{+-},
        \qquad
        \lambda_t^- = s(t)\bar{\lambda} + Y_t^{-+} + Y_t^{--}.
    \]
    For \(i,j\in\{+,-\}\):
    \begin{gather*}
        dY_t^{ij} = -\beta Y_t^{ij}\,dt + \alpha_{ij}\,dN_t^{j} + \rho_{ij}\,dJ_t.\\
        Y_0^{ij}\ge 0,\qquad i,j\in\{+,-\}.\\
    \end{gather*}
    \(N_t^\pm\) are counting processes with intensities \(\lambda_t^\pm\).
    Seasonality modulates only the immigrant baseline \( \bar{\lambda} \), while
    self- and cross-excitation remain in the common reaction state \(Y_t\).

    \section*{4. Deterministic Seasonality and News Intensity}
    \[
        s(t)=\exp\!\Big(
        f_{\mathrm{intra}}(\mathrm{tod}(t))
        +f_{\mathrm{dow}}(\mathrm{dow}(t))
        +f_{\mathrm{doy}}(\mathrm{doy}(t))
        +f_{\mathrm{special}}(t)
        \Big).
    \]

    Let news arrivals be a superposition
    \[
        N_t^{\mathrm{news}} = N_t^{\mathrm{known}} + N_t^{\mathrm{surprise}},
    \]
    where \(N_t^{\mathrm{known}}\) tracks known-in-advance scheduled times and
    \(N_t^{\mathrm{surprise}}\) captures unscheduled shocks.

    On an implementation interval \(h=[t_h,t_{h+1})\), known arrivals are the
    persisted scheduled timestamps in that interval.
    The known component's interval intensity is therefore the scheduled count divided by the interval duration:
    \[
        \nu_h^{\mathrm{known}}
        =
        \frac{
            \left|\{\tau_m^{\mathrm{known}}\in[t_h,t_{h+1})\}\right|
        }{\Delta_h}.
    \]
    The surprise component uses a persisted baseline intensity and a
    non-negative deterministic interval multiplier \(m_h^{\mathrm{surprise}}\):
    \[
        \nu_h^{\mathrm{surprise}}
        =
        \nu_0^{\mathrm{surprise}}m_h^{\mathrm{surprise}}.
    \]
    Equivalently, the deterministic calendar profiles below materialize the
    interval multipliers consumed by the runtime:
    \begin{gather*}
        \nu_{\mathrm{news}}(t)=\nu_{\mathrm{known}}(t)+\nu_{\mathrm{surprise}}(t),\\
        \nu_{\mathrm{surprise}}(t)
        =
        \exp\!\Big(
        g_{\mathrm{dow}}(\mathrm{dow}(t))
        +g_{\mathrm{doy}}(\mathrm{doy}(t))
        +g_{\mathrm{holiday}}(t)
        +g_{\mathrm{session}}(\mathrm{session}(t))
        \Big).\\
    \end{gather*}
    Arrivals are allowed in both open and closed windows, while direct gap increments remain closed-only in \(G_t\).
    Trade seasonality \(s(t)\) modulates the baseline trade-arrival intensity only and does not modulate \(\nu_{\mathrm{news}}(t)\).

    \section*{5. News Marks, Activity Shock, and Gap Process}
    At news times \(\tau_m^{c}\), with \(c\in\{\mathrm{known},\mathrm{surprise}\}\):
    \[
        U_m\sim\Gamma(k_U,\theta_U),\qquad
        \varepsilon_m\sim\text{Kou}(p,\eta_+,\eta_-).
    \]
    Here \(\Gamma(k_U,\theta_U)\) is the shape-scale parameterization (shape \(k_U\), scale \(\theta_U\)).

    Kou base-shock density:
    \[
        f_{\varepsilon}(x)=
        \begin{cases}
            p\,\eta_+e^{-\eta_+ x}, & x\ge 0,\\
            (1-p)\,\eta_-e^{\eta_- x}, & x<0,
        \end{cases}
        \qquad
        \eta_+,\eta_->0,\quad p\in(0,1).
    \]

    The shared severity and return mark are component-scaled before dispatch.
    Let \(a_c\ge0\) denote the activity-coupling scale and \(b_c\ge0\) the
    return-jump scale for component \(c\).
    Then
    \begin{gather*}
        Z_m^{c}=b_c\sigma_G\sqrt{U_m^{c}}\,\varepsilon_m^{c},\\
        J_t^{\mathrm{known}}=\sum_{\tau_m^{\mathrm{known}}\le t} a_{\mathrm{known}}U_m^{\mathrm{known}},
        \qquad
        J_t^{\mathrm{surprise}}=\sum_{\tau_m^{\mathrm{surprise}}\le t} a_{\mathrm{surprise}}U_m^{\mathrm{surprise}},
        \qquad
        J_t=J_t^{\mathrm{known}}+J_t^{\mathrm{surprise}}.\\
        M_t^{\mathrm{known}}=\sum_{\tau_m^{\mathrm{known}}\le t}\mathbf{1}_{\mathrm{open}}(\tau_m^{\mathrm{known}})\,Z_m^{\mathrm{known}},
        \qquad
        M_t^{\mathrm{surprise}}=\sum_{\tau_m^{\mathrm{surprise}}\le t}\mathbf{1}_{\mathrm{open}}(\tau_m^{\mathrm{surprise}})\,Z_m^{\mathrm{surprise}},\\
        M_t=M_t^{\mathrm{known}}+M_t^{\mathrm{surprise}}.\\
        G_t^{\mathrm{known}}=\sum_{\tau_m^{\mathrm{known}}\le t}\mathbf{1}_{\mathrm{closed}}(\tau_m^{\mathrm{known}})\,Z_m^{\mathrm{known}},
        \qquad
        G_t^{\mathrm{surprise}}=\sum_{\tau_m^{\mathrm{surprise}}\le t}\mathbf{1}_{\mathrm{closed}}(\tau_m^{\mathrm{surprise}})\,Z_m^{\mathrm{surprise}},\\
        G_t=G_t^{\mathrm{known}}+G_t^{\mathrm{surprise}}.\\
    \end{gather*}
    Open-window news jumps are applied on the first post-arrival open tick via \(M_t\).
    Closed-window news jumps are applied on the first reopening tick via \(G_t\).

    \section*{6. NIG Business-Time Drivers}
    For \(L^\pm\), define exponents \(\psi_\pm\) by
    \[
        \mathbb{E}\!\left[e^{iuL_t^\pm}\right]=\exp\!\left(t\,\psi_\pm(u)\right).
    \]
    With shared NIG core and one skew shift:
    \begin{gather*}
        \psi_+(u)=\psi_{\mathrm{NIG}}(u;\alpha,\beta_0,\delta,\mu),\\
        \psi_-(u)=\psi_{\mathrm{NIG}}(u;\alpha,\beta_0+\Delta_\beta,\delta,\mu),
        \qquad
        \Delta_\beta<0.\\
    \end{gather*}
    Let the side mean per unit business time be
    \[
        m_\pm=\mu+\delta\frac{\beta_\pm}{\sqrt{\alpha^2-\beta_\pm^2}},
        \qquad
        \beta_+=\beta_0,\quad \beta_-=\beta_0+\Delta_\beta.
    \]
    Trading-time increments use the buy-minus-sell convention with side mean compensation.
    Decompose integrated side activity into deterministic baseline opportunity and endogenous
    Hawkes excitation activity:
    \[
        A_h^\pm = B_h^\pm + E_h^\pm,
    \]
    where \(B_h^\pm\) is generated by the seasonal baseline intensity \(s_h\bar{\lambda}\),
    and \(E_h^\pm\) is generated by the current excitation state and its descendants.
    The NIG return clocks are
    \[
        A_{\mathrm{return},h}^+
        =
        B_h^+ + \eta_+E_h^+,
        \qquad
        A_{\mathrm{return},h}^-
        =
        B_h^- + \eta_-E_h^-,
    \]
    with \(\eta_+>0\) and \(\eta_->0\).
    Calendar seasonality controls baseline event opportunity.
    Hawkes excitation controls stochastic volatility through
    realized endogenous activity.
    Unequal buy/sell excitation return-clock exposure controls skew.
    The simulated trading-time increment is
    \[
        X_{\mathrm{trading}}
        =
        (L^+_{A_{\mathrm{return}}^+}-A_{\mathrm{return}}^+m_+)
        -
        (L^-_{A_{\mathrm{return}}^-}-A_{\mathrm{return}}^-m_-).
    \]
    Thus, the trading-time block has zero conditional drift after conditioning on
    decomposed activity:
    \[
        \mathbb{E}\!\left[
            X_{\mathrm{trading}}
            \mid B^+,E^+,B^-,E^-
        \right]=0.
    \]
    Conditional characteristic function of \(\Delta X_t=X_{t+\Delta}-X_t\):
    \[
        \mathbb{E}\!\left[e^{iu\Delta X_t}\mid B^+,E^+,B^-,E^-\right]
        =
        \exp\!\Big(
            A_{\mathrm{return}}^+(\psi_+(u)-ium_+)
            +
            A_{\mathrm{return}}^-(\psi_-(-u)+ium_-)
        \Big)\,\varphi_{\Delta M+\Delta G}(u).
    \]

    \section*{7. Structured Asymmetry and Constraints}
    \begin{gather*}
        \alpha_{++}=a_{\mathrm{self}},
        \quad
        \alpha_{--}=a_{\mathrm{self}}+\Delta_{\mathrm{self}},
        \quad
        \alpha_{+-}=a_{\mathrm{cross}},
        \quad
        \alpha_{-+}=a_{\mathrm{cross}}+\Delta_{\mathrm{cross}}.\\
        \rho_{++}=r_{\mathrm{self}},
        \quad
        \rho_{--}=r_{\mathrm{self}}+\Delta^{(r)}_{\mathrm{self}},
        \quad
        \rho_{+-}=r_{\mathrm{cross}},
        \quad
        \rho_{-+}=r_{\mathrm{cross}}+\Delta^{(r)}_{\mathrm{cross}}.\\
        \bar{\lambda}>0,\quad
        \beta>0,\quad
        \alpha_{ij}\ge 0,\quad
        \rho_{ij}\ge 0.\\
    \end{gather*}
    With
    \[
        \mathbf{A}_{\alpha}=
        \begin{pmatrix}
            \alpha_{++} & \alpha_{+-}\\
            \alpha_{-+} & \alpha_{--}
        \end{pmatrix},
    \]
    impose Hawkes stability:
    \[
        \rho\!\left(\mathbf{A}_{\alpha}/\beta\right)<1,
    \]
    where \(\rho(\cdot)\) denotes spectral radius.
    This is the excitation branching condition.
    It is independent of deterministic baseline seasonality because \(s(t)\) does not scale descendant excitation.

    \section*{8. Layer Independence Assumption}
    \[
        \mathcal{L}\ \perp\!\!\!\perp\ \mathcal{H}\ \perp\!\!\!\perp\ \mathcal{N},
    \]
    where \(\mathcal{L}\) is the Lévy layer \((L^+,L^-)\), \(\mathcal{H}\) is the Hawkes-layer primitive innovation source for \(N^\pm\), and \(\mathcal{N}\) is the news layer \(\big(N^{\mathrm{news}},\{U_m,\varepsilon_m\}_{m\ge1}\big)\).
    No additional cross-layer dependence is introduced beyond the explicit coupling through \(dJ_t\) in \(dY_t^{ij}\).

    \section*{9. Affine One-Interval Transform}
    For piecewise-constant \(s_h,\nu_h\) on \([0,\Delta]\), with \(\nu_h=\nu_h^{\mathrm{known}}+\nu_h^{\mathrm{surprise}}\), define
    \[
        F(u_+^B,u_+^E,u_-^B,u_-^E,z_+,z_-\mid \mathbf{Y}_0)
        =
        \mathbb{E}\!\left[
                        e^{-u_+^B B^+-u_+^E E^+-u_-^B B^--u_-^E E^-}\,
                        z_+^{N^+}z_-^{N^-}
                        \mid \mathbf{Y}_0
        \right]
        =
        \exp\!\big(C(\Delta)+\mathbf{D}(\Delta)\cdot\mathbf{y}_0\big),
    \]
    with \(C(0)=0,\ \mathbf{D}(0)=\mathbf{0}\), and
    \begin{gather*}
        \mathbf{y}_t=(Y_t^{++},Y_t^{+-},Y_t^{-+},Y_t^{--}),
        \quad
        \mathbf{k}_+=(\alpha_{++},0,\alpha_{-+},0),
        \quad
        \mathbf{k}_-=(0,\alpha_{+-},0,\alpha_{--}),\\
        \mathbf{r}=(\rho_{++},\rho_{+-},\rho_{-+},\rho_{--}),\\
        H_+^E(\mathbf{D})=z_+\exp(\mathbf{D}\cdot\mathbf{k}_+)-1-u_+^E,
        \quad
        H_-^E(\mathbf{D})=z_-\exp(\mathbf{D}\cdot\mathbf{k}_-)-1-u_-^E.\\
        \begin{aligned}
            D_{++}'(\tau) &= -\beta D_{++}(\tau) + H_+^E(\mathbf{D}(\tau)),\\
            D_{+-}'(\tau) &= -\beta D_{+-}(\tau) + H_+^E(\mathbf{D}(\tau)),\\
            D_{-+}'(\tau) &= -\beta D_{-+}(\tau) + H_-^E(\mathbf{D}(\tau)),\\
            D_{--}'(\tau) &= -\beta D_{--}(\tau) + H_-^E(\mathbf{D}(\tau)),
        \end{aligned}\\
        C'(\tau)
        =
        s_h\bar{\lambda}\big(
            H_+^E(\mathbf{D}(\tau))+u_+^E-u_+^B
            +
            H_-^E(\mathbf{D}(\tau))+u_-^E-u_-^B
        \big)
        +\nu_h\Big(\big(1-\theta_U(\mathbf{D}(\tau)\cdot\mathbf{r})\big)^{-k_U}-1\Big),\\
    \end{gather*}
    where \(\nu_h=\nu_h^{\mathrm{known}}+\nu_h^{\mathrm{surprise}}\).
    Setting \(u_\pm^B=u_\pm^E\) recovers the total integrated-activity query.
    Pricing and sufficient-statistics paths can therefore request either total
    activity or separate baseline and excitation loadings from the same affine
    oracle.

    The domain condition is
    \[
        \mathbf{D}(\tau)\cdot\mathbf{r}<\frac{1}{\theta_U}\quad\forall\tau\in[0,\Delta].
    \]

    \section*{10. Anchored Local Martingale Equation}
    Under pricing, the augmented state is
    \[
        Z_h=(\Sigma_h,p_h^Q,D_h^Q).
    \]

    The risk-neutral local characteristic transform is
    \[
        \psi_h^Q(u \mid Z_h)
        =
        \psi_h^P(u+i\theta_h(Z_h)\mid \Sigma_h)-\psi_h^P(i\theta_h(Z_h)\mid \Sigma_h).
    \]
    Equivalently, for the one-interval log-moment-generating function,
    \[
        K_h^Q(v \mid Z_h)
        =
        K_h^P(v-\theta_h(Z_h)\mid \Sigma_h)-K_h^P(-\theta_h(Z_h)\mid \Sigma_h).
    \]

    Define the risk-neutral expected core increment by
    \[
        \mu_h^Q(Z_h)
        =
        \partial_u K_h^Q(u \mid Z_h)\big|_{u=0}
        =
        \mathbb E^Q[R_h \mid Z_h].
    \]

    Then the risk-neutral pressure and anchor recursions are
    \begin{gather*}
        p_{h+1}^Q = \psi_h p_h^Q + (1-\psi_h)\mu_h^Q(Z_h),\\
        A_{h+1}^Q = A_h^Q + c\big(p_{h+1}^Q - p_h^Q\big),\\
    \end{gather*}
    and the risk-neutral deviation recursion is
    \[
        D_{h+1}^Q
        =
        e^{-\kappa \Delta_h} D_h^Q + R_h + A_h^Q - A_{h+1}^Q
        =
        e^{-\kappa \Delta_h} D_h^Q + R_h - c\big(p_{h+1}^Q - p_h^Q\big),
    \]
    and the one-interval observed return satisfies the local martingale condition
    \[
        K_h^Q(1 \mid Z_h)+\big(e^{-\kappa\Delta_h}-1\big)D_h^Q=(r-q)\Delta_h.
    \]

    The pricing engine therefore generates the risk-neutral pressure state and
    anchor from the same local tilted transform used to price.

    \section*{11. SufficientStatistics-Neutral Esscher Tilt}
    \[
        \psi_Q(u)=\psi_P(u+i\theta)-\psi_P(i\theta),
        \qquad
        \mathbb{E}^Q\!\left[e^{X_{t+\Delta}-X_t}\mid\mathcal{F}_t\right]=e^{(r-q)\Delta}.
    \]
    In the anchored setting, the scalar Esscher parameter becomes the local
    state-dependent tilt \(\theta_h(Z_h)\) defined above.

    \section*{12. Final Price Process}
    \[
        X_t=\log\!\left(\frac{S_t}{S_0}\right),
        \qquad
        S_t=S_0\,e^{X_t}.
    \]
    When an observed price tick size \(\delta_S>0\) is supplied or inferred, the
    simulation observation layer projects latent prices to the tick grid by
    nearest-tick rounding from the previous observed price anchor:
    \[
        S_t^{\mathrm{obs}}
        =
        S_{t^-}^{\mathrm{obs}}
        +
        \delta_S
        \operatorname{round}\!\left(
            \frac{S_t^{\mathrm{latent}}-S_{t^-}^{\mathrm{obs}}}{\delta_S}
        \right),
        \qquad
        X_t^{\mathrm{obs}}=\log(S_t^{\mathrm{obs}}/S_0).
    \]
    If no tick size is available, the observation layer uses the latent log-price
    path directly.
    \begin{gather*}
        dS_t
        =
        S_{t-}\,dX_t
        +\frac{1}{2}S_{t-}\,d[X^c]_t
        +S_{t-}\!\int_{\mathbb{R}}\!\big(e^{z}-1-z\big)\,\mu^X(dt,dz),\\
        \Delta S_t = S_{t-}\big(e^{\Delta X_t}-1\big).\\
    \end{gather*}

\end{document}
